\section{When the original polar bracket is \texorpdfstring{\(\F_4\)}{F4}-bilinear}
\label{sec:compatible}

Return to the normal form \eqref{eq:normal-form} and keep the intrinsic
scalar actions on \(U=V/W\) and \(W\). Its auxiliary bracket
\eqref{eq:auxiliary-bracket} is always \(\F_4\)-bilinear.
The following definition concerns the original polar bracket of \(f\).

\begin{definition}\label{def:compatible}
A metabelian solution is \emph{compatible} if \(V\) admits an
\(\F_4\)-vector space structure that restricts to the intrinsic structure
on \(W\), induces the intrinsic structure on \(U\), and makes the
original polar bracket \(\F_4\)-bilinear.
\end{definition}

\begin{theorem}\label{thm:compatibility}
For a normal form with data \(P_i,a_{ij}\), the following are equivalent:
\begin{enumerate}
\item the solution is compatible;
\item there are \(b_1,\ldots,b_d\in W\) satisfying
\begin{equation}\label{eq:compatibility-system}
 P_ib_i=0,\qquad P_ib_j=P_ja_{ij}\quad(i\ne j);
\end{equation}
\item the auxiliary \(\F_4\)-Lie algebra has a restricted square \(r\)
with \(r(W)=0\) and \(r(\mathfrak l)\subseteq W\).
\end{enumerate}
If a compatible scalar lift exists, the number of such scalar actions
on the fixed binary space, with the prescribed actions on \(U,W\), is
\(4^{dk}\), where \(k=\dim_{\F_4}\bigcap_i\ker P_i\).
\end{theorem}
\begin{proof}
Every binary-linear lift of multiplication by \(\omega\) has the form
\[
 T(x,z)=(\omega x,\omega z+t(x)).
\]
The relation \(T^2+T=\id\) is exactly
\(t(\omega x)=\omega^2t(x)\). Thus
\(t(x)=b(x)=\sum_i x_i^2b_i\).
The involutory binary shear \(S(x,z)=(x,z+b(x))\) conjugates \(T\)
to standard scalar multiplication. It changes the cocycle of \(f\) to
\[
 Q+b+P(Q+b).
\]
The first two terms have \(\F_4\)-bilinear polars. The last term has
scalar weight zero. Its polar could be \(\F_4\)-bilinear only if that
polar were zero, since an \(\F_4\)-bilinear polar has weight two.
An additive zero-constant function of weight zero is zero: additivity
and invariance under \(\omega\) give its value at \(\omega^2x=x+\omega x\)
equal to zero, and then invariance gives its value at \(x\) equal to zero.
Compatibility is therefore equivalent to
\begin{equation}\label{eq:compatibility-polynomial}
 P(x)(Q(x)+b(x))=0.
\end{equation}
After the triangle cancellations, comparison of \(x_i^3\) and
\(x_ix_j^2\) gives precisely \eqref{eq:compatibility-system}.

We recall explicitly the restricted-square axioms in characteristic two:
\(r(\alpha u)=\alpha^2r(u)\),
\(r(u+v)=r(u)+r(v)+[u,v]_0\), and
\(\ad r(u)=(\ad u)^2\).
Any square with the range and vanishing conditions in (iii) must have
the form
\begin{equation}\label{eq:restricted-square}
 r(x,z)=(0,P(x)z+q(x)),\qquad q(x)=Q(x)+b(x).
\end{equation}
Its first two axioms are automatic. For \(u=(x,z),v=(y,t)\), the two
adjoint expressions have final coordinates
\[
 P(x)P(y)z+P(x)\beta(x,y),\qquad
 P(x)P(y)z+P(y)q(x),
\]
respectively. Condition~\eqref{eq:compatibility-polynomial} implies
\(P(x)\beta(x,y)=P(y)q(x)\): take the coefficient of \(s\) in
\(P(x+sy)q(x+sy)=0\). Its degree in \(s\) is at most three, below
\(|\F_4|\), so the comparison is valid for functions over \(\F_4\).
This proves the third restricted axiom. Conversely that axiom, applied
to \(u\), gives \([u,r(u)]_0=0\), hence
\eqref{eq:compatibility-polynomial}.

Finally two solutions of \eqref{eq:compatibility-system} differ by
independent choices of \(d\) vectors in \(\bigcap_i\ker P_i\).
The displayed count follows. It counts scalar actions, not their
conjugacy classes.
\end{proof}

\begin{example}\label{ex:incompatible-eight}
The eight-dimensional class-three model has
\(P_1z=y\), \(P_1y=0\), \(P_2=0\), and \(a_{12}=z\).
Condition~\eqref{eq:compatibility-system} would require
\(0=P_2b_1=P_1a_{12}=y\). Hence it has no compatible scalar action.
This gives a concrete distinction between the auxiliary algebra and
the original polar algebra.
\end{example}

\begin{remark}
The restricted square in Theorem~\ref{thm:compatibility} is a square on
the auxiliary algebra. The theorem imposes both its range and its
restriction to the derived ideal. General correspondences between
quadratic maps and restricted Lie structures, such as
\cite[Section~9]{Pakianathan2007}, carry their own additional identities.
Here the exact additional condition is the solvable linear system
\eqref{eq:compatibility-system}.
\end{remark}
